% Source: 43_Chapter_32_Aiming_At_the_Foundations_The_Cult_of_R.txt
\bookchapter[Aiming at the Foundations]{CHAPTER 32}{AIMING AT THE FOUNDATIONS}{The Cult of Reason}

\leadin{I}{F AMERICANS BELIEVED IN} saints, Benjamin Franklin would be among them. He exemplified so
many virtues Americans have come to admire. People found him practical, earthy, affable, witty and,
above all, tolerant.

A few weeks before he died, Ben responded to an inquiry by President Ezra Stiles of Yale concerning
his religious faith. Said Franklin,

\begin{blockquote}
As to Jesus of Nazareth~.~.~. I have~.~.~. some doubts as to his Divinity, tho’ it is a question I
do not dogmatize upon, having never studied it, and think it needless to busy myself with it now,
when I expect soon an opportunity of knowing the truth with less trouble. I see no harm, however, in
its being believed, if that belief has the good consequence~.~.~. of making his doctrines more
respected and better observed.
\end{blockquote}

Something of the American spirit is there. It is the spirit of Franklin’s time, the Age of Reason
(1648–1789). Questions of dogma seemed unimportant, hardly worth fretting about. What was immensely
more important was behavior. Do our beliefs make us more tolerant, more respectful of those who
differ with us, more responsive to the true spirit of Jesus?

If that hatred of religious bigotry, coupled with a devotion to tolerance of all religious opinions,
has a familiar ring, it is because the attitudes of the Age of Reason are not a thing of the past.
They live on today in the values of the Western world.

The Age of the Reformation proved again that faith and power are a potent brew. As long as
Christians had access to power, they used it to compel conformity to the truth: Catholic, Lutheran,
or Reformed. So men died for their faith, tens of thousands of them, until something general but
very deep in man awoke to revolt.

\section{THE SPIRIT OF THE AGE OF REASON}

We call that revolt the Age of Reason—or as some prefer, the Enlightenment—when, as Charles Williams
put it, “national interests and mental relaxations combined to exclude metaphysics from culture.”

The spirit of the Age of Reason was nothing less than an intellectual revolution, a whole new way of
looking at God, the world, and oneself. It was the birth of secularism.

The Middle Ages and the Reformation were centuries of faith in the sense that reason served faith,
the mind obeyed authority. To a Catholic it was church authority, to a Protestant biblical
authority, but in either case God’s Word came first, not man’s thoughts. Man’s basic concern in this
life was his preparation for the next.

The Age of Reason rejected that. In place of faith it set reason. Man’s primary concern was not the
next life, but happiness and fulfillment in this world; and the mind of man, rather than faith, was
the best guide to happiness, not emotions, or myths, or superstitions.

The spirit and purpose of the Enlightenment were eloquently expressed by one of its spokesmen, Baron
von Holbach, who wrote,

\begin{blockquote}
Let us endeavor to disperse those clouds of ignorance, those mists of darkness, which impede Man on
his journey,~.~.~. which prevent his marching through life with a firm and steady step. Let us try
to inspire him~.~.~. with respect for his own reason—with an inextinguishable love of truth~.~.~. so
that he may learn to know himself~.~.~. and no longer be duped by an imagination that has been led
astray by authority~.~.~. so that he may learn to base his morals on his own nature, on his own
wants, on the real advantage of society~.~.~. so that he may learn to pursue his true happiness, by
promoting that of others~.~.~. in short, so that he may become a virtuous and rational being, who
cannot fail to become happy.
\end{blockquote}

\section{THE RENAISSANCE}

How did this new spirit take root and grow? Its seeds probably lie in the Reformation era itself, in
a movement historians call the Renaissance. The word means rebirth and refers to the recovery of the
values of classical Greek and Roman civilization expressed in literature, politics, and the arts. In
one sense, the Reformation is not possible without the influence of Renaissance humanism, which
seeks human flourishing, a well-rounded person who possesses a full complement of virtues and
abilities. Reading classic texts from the ancient Greek and Roman glory days was thought to be a
crucial step towards living well. To read was to engage a famous author in a conversation, an
experiential encounter with both author and culture. Reading a text in its original language was the
gold standard in this soul-building venture. Each of the major Reformers embraced this Renaissance
humanist approach to texts. Each came to believe that they must encounter the Bible in its original
languages.

Perhaps the best example of the rebirth of the classical spirit was Erasmus of Rotterdam
(1467–1536). In a series of best-selling satires highlighted by Praise of Folly, Erasmus ridiculed
monasticism and Scholasticism by the use of irony, wit, and enlightened common sense.

Erasmus’s followers thought they heard a true reformer in him. He laid the egg, they said, that
Luther hatched. But in 1524 a significant conflict between Luther and Erasmus erupted. In that year
Erasmus’s Diatribe on Free Will appeared. It made clear the cardinal differences between the two
men. Luther believed that the human will was enslaved, totally unable, apart from an enabling grace,
to love or serve God. But Erasmus considered this a dangerous doctrine since it threatened to
relieve man of his moral responsibility.

The differences in the Reformation and the Renaissance lie right there, in the view of man. The
Reformers preached the original sin of man and looked upon the world as fallen from God’s intended
place. The Renaissance had a positive estimate of human nature and the universe itself. This
confidence in man and his powers flowered and filled the air with fragrance during the
Enlightenment.

Another root of the Enlightenment runs through the century (1550–1650) of appalling religious
conflicts: the English Civil War, the persecution of the French Huguenots, and the Thirty Years’ War
in Germany. Common decency cried out against the power of fanatical clerics. More and more people
felt only disgust at the burning or drowning of an elderly woman accused of witchcraft or heresy.
Religious prejudice seemed like a far greater danger than atheism. So a thirst for tolerance and
truths common to all men spread.

\begin{sidebar}{Enlightenment}
The Enlightenment was actually a series of enlightenments. It influenced England, France, and
elsewhere before finally beginning in Germany. Its basic insight was captured in the first French
Enlightenment thinkers. They made several observations from their experiences in England, which
fascinated them. They saw greater prosperity and greater liberty. In England, the institutions of
government and church did less to restrict individual citizens who were free to reason for
themselves to a greater degree. These several elements were finally gathered into a central
conviction: if individuals were not constrained and restricted by the burden of cultural traditions
and were allowed to freely reason for themselves, they would enjoy prosperity. Freedom to think for
oneself, independent from tradition, yields truth and prosperity. This insight is expanded and
extended in endless varieties.
\end{sidebar}

Finally, the Age of Reason sprang from the soil of a new faith in law and order. Modern science
arose in the sixteenth and seventeenth centuries and filled men with visions of a new day of peace
and harmony. These pioneers of modern science forced men to think in a new way about the universe:
Copernicus (1473–1543), who insisted that the sun, not the earth, was the center of our universe;
Johann Kepler (1571–1630), who concluded that the sun emitted a magnetic force that moved the
planets in their courses; and Galileo Galilei (1564–1642), who made a telescope to examine the
planets and proved that the acceleration of falling bodies is constant.

All these discoveries, however, had to be united in one all-embracing principle that would explain
the motion of bodies in the heavens and present the universe as one great machine operating
according to unalterable laws. This was the feat of the most illustrious scientist of the Age of
Reason, Isaac Newton (1642–1727).

In 1687 Newton published his momentous work, Mathematical Principles of Natural Philosophy, in which
all laws of motion, in the heavens and on the earth, were harmonized in a master principle for the
universe, the law of gravitation.

\section{THE WONDER OF THE WORLD-MACHINE}

The reading public of Europe was captivated by the wonder of Newton’s world-machine. The medieval
world of unseen spirits—angels and demons—seemed less plausible or even superstitious. In its place
was a sun-centered universe that operated by physical laws explained and justified by mathematics.

In the new model, the sun displaced the earth as the center. Some believed that mankind had been
displaced as the crowning apex of creation in the center of God’s world. Some felt that God had been
displaced as well. God seemed less necessary to sustain the world when one observed that the world
operated by discernible laws. Readers today should be reminded that the rise of modern science is
absolutely dependent upon Christian convictions and is carried on by minds trained in Christian
environments and institutions. Also worth noting is the re-envisioning of God. In these discussions,
God is often pictured as one more actor on the stage and not the majestic Lord overseeing the entire
project.

This vast unfolding universe filled some men with dismay. Blaise Pascal (1623–1662), the French
physicist, recoiled before “the terrifying expanses which engulf us as a fleeting atom.” But others
accepted it as an invitation to discover the secrets of cosmic existence. Alexander Pope wrote:

\begin{poem}
Nature and Nature’s laws lay hid in night;\\
God said, “Let Newton be!” and all was light.
\end{poem}

This sudden access to the mysteries of the universe seemed to magnify the role of human reason. If
the universe is a smooth-running machine with all its parts coordinated by one grand design, then
man only has to think clearly to find life’s meaning and true happiness. This fundamental idea, that
man has the ability to find the truth by the use of his senses and reason, gave rise to the label
Age of Reason.

The term should not suggest that every blacksmith and village priest suddenly assumed the airs of an
intellectual. Many Christians lived and died in the faith of their fathers, totally unconscious that
a new age had dawned. But the outlook and direction of Europe had changed.

Christianity could scarcely escape the fallout from this intellectual revolution. For 1,200 years
Augustine’s ideas had ruled Christendom. Man was an enslaved sinner who needed, above all else, the
supernatural grace of God. To insure the availability of this grace through the Christian church,
God had ordained the powers of the state to protect truth and punish error.

But now intellectuals were arguing something else: Man is no sinner. He is a reasonable creature.
Now he seemed to need common sense more than God’s grace.

Christians found themselves in two contrasting climates. At first, during the closing years of the
seventeenth century, some believers, especially in England, tried to harmonize reason and faith.
They argued that Christianity is totally reasonable but some truths come by reason and some by
revelation. Some things, like the existence of God, come by observing the heavens; while other
things, like the resurrection of Christ, come by the witness of Scripture.

In time, however, after the beginning of the eighteenth century, the climate changed. In France
confidence in reason soared and Christians found that many intellectuals dismissed all appeals to
revealed Scripture as superstitious nonsense. The climate was obviously more hostile.

The best representative of the first generation is John Locke (1632–1704). The highly influential
philosopher never minimized the importance of belief. In his Essay Concerning Human Understanding,
he not only shows how reason functions; he indicates that the existence of God is “the most obvious
truth that reason discovers.” Closer examination, however, shows that the God Locke had in mind has
little in common with the God of the Exodus or of Jesus’ resurrection. Attempts to rationally
justify Christianity sadly often led to a revising of Christianity to what seems rational at the
time. A profound ancient Christian wisdom is forgotten. Christians do not arrive at Christian
conviction as the conclusion of a long argument; they believe because they received witness from the
first eye witnesses. This witness was powerfully vindicated by the ongoing power of God. The New
Testament is a book of powerfully vindicated testimony of God’s action, not a book of deductive
logic. Reason may serve revelation, but it can never replace revelation.

To Locke’s credit, he endorses revelation and never doubts its importance. Fulfilled prophecies and
the miracles of Jesus are proof of Jesus’ authority. But what Scripture actually shows is that few
dogmas are necessary. And these are simple and intelligible to ordinary men. Christianity, in fact,
has only one essential doctrine: Jesus is the Messiah. Most of traditional theology Locke casually
dismisses as irrelevant.

Locke also spoke for his generation in his emphasis upon moral conduct. Christianity, he says, adds
to its belief in Jesus as Messiah the imperative of a good life. Jesus often spoke of rewards and
punishments for Christian behavior. And that, too, is entirely reasonable because reason shows that
moral standards must be reinforced by strong incentives.

Thus, according to Locke, revelation shows Christianity’s reasonable character. Belief in Jesus as
Messiah and man’s ethical behavior are all Jesus and the apostles required for righteousness. Both
of these are basically rational.

\section{THE RISE OF THE RADICALS}

“The fathers have eaten sour grapes, and the children’s teeth are set on edge” (Ezek. 18:2). Many in
the next generation, the first of the eighteenth century, felt fewer obligations to the Christian
past, so instead of trying to harmonize nature and Scripture, they simply set aside revelation. Many
intellectuals claimed that the parts of the Bible that agree with reason are clearly unnecessary.
The parts that contradict reason—the myths, miracles, and priestly mumbo jumbo—are simply untrue.
This more militant attitude toward the faith was especially evident in France.

In the eighteenth century Paris arose as the capital of a new cosmopolitan culture. Ideas circulated
freely throughout Europe and the American colonies. To a degree unmatched before or since, the
social and intellectual leaders of Europe were united in a community of thought and interests.

In Paris a group of thinkers and writers known as the philosophes brought the Age of Reason to its
climax. The philosophes were not philosophers devoted to an academic discipline. They were men of
letters, students of society who analyzed its evils and advocated reforms. They aimed to spread
knowledge and emancipate the human spirit.

Curiously enough, atheism was not at all fashionable in this “polite society.” Most of the prominent
“infidels” who ridiculed Christianity during the eighteenth century believed in a supreme being but
regarded it superstitious to hold that he interfered with the world-machine. This belief comes to be
called deism, a movement especially popular among English speakers. Deism served as a halfway house
on the road to atheism. One could keep the idea of God and dismiss the concept that God would engage
or interfere with the world.

The God of the deists has sometimes been called the watchmaker God. God created the world as a
watchmaker makes a watch, and then wound it up and let it run. Since God was a perfect watchmaker,
there was no need of his interfering with the world later. Hence the deists rejected anything that
seemed to be an interference of God with the world, such as miracles or special revelations recorded
in the Bible.

The deists believed that their religion was the original religion of man. From it had come, by
distortion, all other religions. These distortions were the work of priests who concocted the
theologies, myths, and doctrines of the various religions to enhance their own power.

The most influential propagandist for deism was Voltaire (1694–1778), who personified the skepticism
of the French Enlightenment. Above all others, Voltaire popularized Newton’s science, fought for
personal liberty and freedom of the press, and spread the cult of reason. He turned out a prodigious
number of works: histories, plays, pamphlets, essays, and novels. In his correspondence, estimated
at ten thousand letters, he wittily spread the virtues of Enlightenment and scathingly attacked the
abuses of his day.

Voltaire achieved his greatest fame as the most relentless critic of the established churches,
Protestant and Catholic alike. He was sickened by the intolerance of organized Christianity and
disgusted by the petty squabbles that seemed to monopolize the time of many priests and clergymen.
Yet, in spite of his biting criticism of Christianity, his aim was not religion’s destruction. He
once said that if a God did not exist, it would be necessary to invent one.

Voltaire had many disciples, but his only serious rival in spreading the gospel of deism was a set
of books: the famous French Encyclopedia edited by Denis Diderot (1713–1784). The seventeen volumes
of the Encyclopedia constituted the chief monument of the philosophes. They heralded the supremacy
of the new science, championed tolerance, denounced superstition, and expounded the merits of deism.
Diderot’s article on Christianity professed high regard for the religion of Jesus, but its effect
was to stir the reader to a profound contempt for Christianity’s social failures.

Unlike most previous critics of the church, the philosophes were not heretics or dissenters who
attacked the church in the name of Christ. These men launched their attack from outside the church.
And they aimed their missiles not at a single point of dogma but at the foundation of all Christian
truth. Their well-publicized purpose was to demolish the citadel.

\section{CHRISTIANITY ON TRIAL}

Christianity, they insisted, was a pernicious plot designed to turn the earth over to the oppressive
powers of a priestly caste. Revealed religion was nothing less than a scheme to exploit the
ignorant. Voltaire liked to refer to Christianity as the “infamous thing.” His most ruthless charge
against the faith pictured the thousands upon thousands of victims of Christianity’s intolerance.
Ironically, Christianity’s critics measured Christian behavior by standards they called human. These
critics ignored the fact that these human standards were actually the legacy of Christian teaching.

These intellectuals judged Christianity by the simple human standards of good and evil. If the
church, in the name of purity of doctrine, sanctioned the bloody carnage of fellow Christians, as it
had in the wars of religion, then Christianity, far from being sacred and holy, was a wicked
institution. It had prevented peace, harmony, and progress among the peoples of the earth.

The primary weapon aimed at the church was truth. “We think that the greatest service to be done to
men,” said Diderot, “is to teach them to use their reason, only to hold for truth what they have
verified and proved.”

But the standards of truth ruled out Christian doctrine from the start. When the orthodox tried to
reason from their basic premises, the infidels only scoffed because they refused to allow arguments
drawn from authority or tradition embodied in the Bible or the church. These simply were not
“reasonable.”

Appeals to miracles met with similar disdain. The proof of a position was found in reason or human
experiences, and since miracles failed this test they were dismissed as medieval nonsense.

“You see,” Diderot argued,

\begin{blockquote}
once one sets foot in this realm of the supernatural, there are no bounds, one doesn’t know where
one is going nor what one may meet. Someone affirms that five thousand persons have been fed with
five small loaves; this is fine! But tomorrow another will assure you that he fed five thousand
people with one small loaf, and the following day a third will have fed five thousand with the wind.
\end{blockquote}

The critics were thoroughly aware that they were fomenting a revolution in the fundamental beliefs
of Europeans. Voltaire reported on each new triumph of reason over the church with the exultation of
a commander winning battles.

\section{THE FUTILITY OF EMPTY CONFIDENCE}

These attacks upon Christian convictions demanded a vigorous and well-reasoned response from
orthodox Christians. Unfortunately in Catholic countries, like France, the reaction was woefully
inadequate. Church leaders were not indifferent to the rising tide of infidelity, but they tried to
check it by traditional means. They appealed to secular authorities to censor the “dangerous” books,
but they were usually unfamiliar with the primary issues the scoffers raised.

In England it was different. Several men wrote effectively against deism, none more so than Bishop
Joseph Butler (1692–1752). His monumental work, The Analogy of Religion, virtually ended the debate
for thinking people. Skirmishes continued for years, but after Butler it was clear the fundamental
issues had been settled.

The deists, with their confident optimism, assumed that they knew all about God’s wisdom and
purpose. They read it all in the pattern of nature. Butler, however, saw with disarming clarity that
life is filled with perplexities and enigmas.

He did not try to prove the existence of God. Deists never denied this premise. Nor did he reject
reason; he accepted it as man’s natural light. But he did challenge reason’s sovereignty. Reason,
said Butler, provides no complete system of knowledge, and in ordinary life it can offer us only
probabilities.

Thus Butler undermined deism’s fortress, its confidence in reason. Nature, he said, is not a realm
in which reason is supreme. It is filled with obscurities and unexplained mysteries. We meet
perplexity at every turn. And if we meet problems in nature, should we be surprised if we encounter
difficulties in religion?

But take one further step. We know the ordinary course of nature because it is disclosed to us by
experience, including our problems and perplexities. If religious truths encounter similar
difficulties, isn’t it reasonable to assume that one kind of knowledge is as dependable as the
other? We normally act on the basis of probability. Why not in religion?

In the end, however, deism collapsed from its own weaknesses. It was based on a false optimism. It
had no explanation for the evils and disasters of life. Because the laws of nature were clear and
unalterable, deists assumed that man’s moral choices drawn from nature were also simple and
unchanging. If asked, “Why don’t men always see clearly the religious truths in nature?” the deist
could only respond with, “the lies of priestcraft.” But that was too simple to be true, and few were
convinced.

The eventual rejection of deism, however, did not restore Christianity to a central place in Western
culture. The negative work of the Age of Reason endured. Modern culture—its art, its education, its
politics—was freed from formal Christian influence. Men made a deliberate attempt to organize a
religiously neutral civilization. This meant that faith was to function in the private realms
confined to church, home, and heart. That is what we find today in modern secular societies.

This leaves Christians with a basic problem in the modern era: How far can believers go in trying,
as citizens, to get the state to enforce Christian standards of conduct? Or if Christians give up
the idea of enforcing Christian behavior, then what norm of conduct should they, as citizens, try to
make an obligation for everyone?

\readinghead{Suggestions for Further Reading}
\begin{readinglist}
\entry Brinton, Crane. The Shaping of the Modern Mind. New York: The New American Library, 1953.
\entry Cragg, Gerald R. The Church and the Age of Reason: 1648–1789. Middlesex: Penguin Books, 1960.
\entry Gay, Peter. Age of Enlightenment. New York: Time-Life Books, 1966.
\entry ———. The Enlightenment: An Interpretation: The Rise of Modern Paganism. New York: Alfred A.
  Knopf, 1966.
\entry McGrath, Alister. Science and Religion: A New Introduction, 2nd edition. Malden MA:
  Blackwell, 2010.
\entry Pearse, Meic. The Age of Reason: From the Wars of Religion to the French Revolution, AD
  1570–1789, Baker History of the Church. Vol. 5. Grand Rapids: Baker, 2006.
\end{readinglist}
